\chapter{Szakirodalmi háttér}

\section{Hálózatkutatás és közösségszerkezet}

A gráfok társas kapcsolatok modellezésében régóta központi szerepet játszanak. A videojátékos közösségek esetében különösen izgalmas kérdés, hogy a játékosok közötti ismétlődő interakciók mennyiben alkotnak felismerhető struktúrákat, szubközösségeket, hidakat és lokális mintázatokat. Girvan és Newman klasszikus munkája rámutatott arra, hogy a hálózat élei nem pusztán kapcsolódási elemek, hanem közösségszerkezetet leíró, eltérő fontosságú viszonyok is lehetnek \cite{girvannewman}. Fortunato átfogó áttekintése pedig jól összefoglalja, miért központi probléma a közösségdetektálás a modern hálózatelemzésben \cite{fortunato2010}.

A jelen projektben a közösségfogalom nem öncélú. Azért fontos, mert a többször ismétlődő kapcsolatokra épülő játékoscsoportok értelmezhetőbbé teszik a nagyméretű hálózatot, segítenek elkülöníteni a valós kapcsolati magokat a zajtól, és támaszt nyújtanak a later\-i útvonal- és analitikai számításoknak.

\section{Közösségdetektálás és modularitás}

A korai Python oldali populációs klaszterezés a NetworkX \textit{greedy\_modularity\_communities} eljárását használta, amely a Clauset--Newman--Moore-féle mohó modularitásmaximalizálási megközelítésre épül \cite{clauset2004,networkx_greedy}. Ez tudatos kompromisszum volt: nem a leglátványosabb, hanem gyorsan magyarázható és reprodukálható megoldás. A jelenlegi futásidejű gráfállapotot a Rust motor önállóan, a meccs-JSON fájlokból és az SQLite-adatbázisból építi fel; Python ebben a szerepben lassú és nehezen skálázható lett volna. A Python pipeline szerepe ezért legacy/populációs közösségek és kompatibilitási exportok előállítása.

A modularitás-optimalizáció ismert korlátait sem hagyhatók figyelmen kívül: Fortunato és Barthelemy eredménye szerint a felbontási korlát miatt kisebb, de valós szerkezetek a globális optimum érdekében összeolvadhatnak \cite{fortunato_barthelemy2007}. A \textit{python\_population} klasztereket ezért nem szabad ontológiai "igazi közösségként" kezelni, inkább a választott gráf, súlymodell és heurisztika melletti legjobb, reprodukálható populációs partícióként.

\section{Asszortativitás}

Newman munkája az asszortatív keveredést olyan hálózati jelenségként írja le, amikor a kapcsolódó csomópontok hasonló vagy éppen eltérő tulajdonságúak \cite{newman2003}. A jelen projekt ezt a játékosok teljesítménymutatóira alkalmazza: nem az a kérdés, hogy a magas fokszámú csúcsok kapcsolódnak-e egymáshoz, hanem az, hogy a gráfban összekapcsolt játékosok hasonló \textit{opscore} vagy \textit{feedscore} értékeket hordoznak-e.

A dolgozat két gráfmódot különböztet meg: a \textit{social-path} csak a pozitív, szövetségesi támaszú kapcsolatokra épít, a \textit{battle-path} az összes ismétlődő együtt- vagy egymás ellen játszást figyelembe veszi. Az asszortativitás mértéke e két módban eltérhet, ami önmagában is informatív.

\section{Esport-analitika és játékos-teljesítménymérés}

A League of Legends és az általánosabb MOBA-játékok esport-analitikája önálló kutatási területté vált. Mora-Cantallops és Sicilia átfogó irodalmi áttekintése megmutatja, hogy a MOBA-játékok egyedi kombinációt kínálnak verseny, együttműködés, játékosélmény és társas interakció között, és ez a kombináció különlegesen gazdag adatforrást teremt empirikus vizsgálatokhoz \cite{moracantallops2019}.

A játékos-teljesítmény modellezése Novak és munkatársai munkájában is megjelenik: profi League of Legends meccsekre kidolgozott teljesítménymodellek rámutatnak, hogy a domainspecifikus változók és a strukturált adatfeldolgozás fontosabb, mint az általános statisztikai modellek \cite{novak2022}. Ez erősíti a jelen projekt szerepkörérzékeny pontozási megközelítésének indokoltságát.

Bahrololloomi és munkatársai a hasonló adatgyűjtési stratégiát (kiindulópontból való mintavétel és ismétlődő meccskapcsolatok felderítése) gépi tanulásos keretben vizsgálják \cite{bahrololloomi2021}.

A csapathálózat és teljesítmény összefüggése Sánchez-González és González-Bailón munkájában is megjelenik, ahol profi LoL csapatok hálózati struktúráját kapcsolják össze a csapathatékonysággal \cite{sanchezgonzalezBailon2023}. Ez megalapozza azt a nézetet, hogy a jelen projektben használt klaszter- és kapcsolatstruktúra nem puszta vizualizáció, hanem potenciálisan prediktív erejű analitikai egység.

Mindezek alapján a jelen projekt a meglévő esport-analitikai irodalom egy természetes empirikus kiterjesztése: nem profi, hanem amatőr/félprofesszionális rangsorolt játékos-hálózatokat vizsgál, és nem egyetlen meccsre, hanem ismétlődő kapcsolatokra épít.

\section{Legrövidebb utak és heurisztikus keresés}

A gráfok másik természetes vizsgálati iránya az útvonalak keresése. A klasszikus legrövidebb út problémát Dijkstra alapozta meg \cite{dijkstra1959}, míg az A* keresés formális alapját Hart, Nilsson és Raphael adták meg \cite{hart1968}. A jelen projektben azért vált ez fontos kérdéssé, mert a játékosháló nemcsak statikus objektumként érdekes, hanem úgy is, mint egy olyan kapcsolati tér, ahol két játékos között különböző ismeretségi láncok húzódhatnak.

A Rust rétegben az egzakt A* implementáció ALT-szerű alsó becsléseket használ, amelyeket a Goldberg--Harrelson-féle megközelítés ihletett \cite{goldberg2005}. A cél nem egy új algoritmus feltalálása, hanem a játékosháló nagyságrendjéhez illeszkedő, helyes és stabil futtatókörnyezet.

\section{Ágensorientált modellezés}

Az ágensorientált modellezés (ABM) olyan szimulációs paradigma, amelyben az egyéni szereplők lokális szabályok alapján cselekszenek, és a makroszintű minták, mint terjeszkedés, verseny vagy egyensúly, kizárólag ezeknek az egyéni döntéseknek az összegzéseként jelennek meg. Bonabeau és munkatársai ezt a keretrendszert természetes és mesterséges rendszerekre egyaránt alkalmazták \cite{bonabeau2002}.

A dolgozat utolsó részét alkotó Tribal NeuroSim kísérlet ebbe a paradigmába illeszkedik. A PremadeGraph pipeline által generált klaszterprofilok törzsi entitásokként jelennek meg azonos szimulációs szabályrendszerben, és a cél annak megfigyelése, hogy a két adatkészlet (Flex Queue és SoloQ) eltérő gráfstruktúrája különböző emergáló viselkedési stratégiákat eredményez-e. Az ABM irodalom éppen azért releváns háttér, mert hangsúlyozza: a szimulációból kapott eredmény a modell viselkedéseként értelmezendő, nem a valóság közvetlen leképezéseként.
